\section{A determinant estimate with marks on one side}
\label{det:section}
\label{sec:determinant}

The first correlation estimate removes prime conditions from factors of
the unmarked endpoint.  Throughout this section and the next two sections,
put
\begin{equation}\label{det:parameters}
 L=\log x,\qquad W=\exp(L^{0.24}),\qquad
 V(W)=\prod_{p\leq W}(1-p^{-1}).
\end{equation}
An integer is \emph{rough} if it has no prime factor at most $W$.
Fix disjoint groups consisting of all the primes in the bands
\begin{equation}\label{det:bands}
 \mathcal P_i=\{p:\exp(L^{a_i})\leq p\leq\exp(2L^{a_i})\},
 \qquad 0.1<a_1<\cdots<a_K<0.2.
\end{equation}
Here $K$ and the $a_i$ are fixed before $x$ tends to infinity.  Write
\begin{equation}\label{det:marks}
 \begin{aligned}
 V_i&=\sum_{p\in\mathcal P_i}\frac1p,
 &\mu_i(p)&=\frac1{pV_i},\\
 \omega_i(h)&=\sum_{p\in\mathcal P_i}\1_{p\mid h},
 &\omega(h)&=\sum_{i=1}^K\omega_i(h).
 \end{aligned}
\end{equation}
Lemma~\ref{lem:prime-estimates} and partial summation give
$V_i=\log 2+o(1)$.  For a fixed $q\in(0,1)$, the marked weight is
\begin{equation}\label{det:weight}
 \mathcal W(h)=q^{\omega(h)-K}
       \prod_{i=1}^K\frac{\omega_i(h)}{V_i}.
\end{equation}
The weight is zero if some group supplies no divisor, and is bounded by
a constant depending only on $K,q$.  Indeed, $q^{j-1}j$ is bounded for
positive integers $j$.

\begin{theorem}[One-sided marked Type II estimate]\label{thm:type-II}
Fix $\delta,C,D_*>0$ and $q\in(0,1)$.  Suppose
$H_m,H_n\geq x^\delta$ and $X=H_mH_n\asymp x$.  Let $F$ satisfy
$|F(h)|\leq1$ and $F(ph)=F(h)$ for every prime in the groups
\eqref{det:bands}.  Let $\alpha_m$ be supported on an arbitrary
interval in $[H_m,2H_m]$, where it has the value
\[
 \alpha_m=m^{iv}\left(\1_{m\text{ prime}}
       -\frac{\1_{m\text{ rough}}}{V(W)\log m}\right),
 \qquad |v|\leq L^C.
\]
Let $\beta_n$ be supported on rough integers in $[H_n,2H_n]$, with
$|\beta_n|\leq L^C$.  If $K$ is sufficiently large in terms of
$\delta,C,D_*,q$, then
\begin{equation}\label{eq:type-II}
 \left|\sum_{m,n}\alpha_m\beta_n F(mn-1)\mathcal W(mn-1)\right|
       \ll X L^{-D_*}.
\end{equation}
The required lower bound on $K$ is independent of the particular
$a_i$.  The implicit constant and the threshold for $x$ may depend on
all the fixed parameters, including the $a_i$.
\end{theorem}

We first give the reduction and geometric construction.  The signed
moment estimate is proved in Proposition~\ref{prop:minor-moment};
Proposition~\ref{prop:major-term} estimates the resulting major term
and completes the proof of the theorem.  Exponents denoted $O_C(1)$
below can be chosen independently of $K$.  A factor depending on a
fixed $K$ is harmless, but a loss $L^{cK}$ would not be harmless at
the final choice of $K$; we keep this distinction explicit.

\subsection{Cauchy's inequality and the normalization of the square}

Expand the product of the $\omega_i$ by selecting one prime from
each group, and denote their product by $b_0$.  For $mn-1=b_0h$,
invariance gives $F(mn-1)=F(h)$.  Except when some selected prime
has square dividing $mn-1$, one also has
$\omega(h)=\omega(mn-1)-K$.  The replacement of the original
damping by $q^{\omega(h)}$ costs
\begin{equation}\label{det:squareful-error}
 O_A(XL^{-A})\qquad\text{for every fixed }A>0.
\end{equation}
To verify this, the sum over labels of either nonnegative weight is
$O_{K,q}(1)$: removing $K$ distinct prime factors can lower $\omega$
by at most $K$, so the new weight is at most the old one.
For fixed $m$ and a group prime $p$, the congruence
$mn\equiv1\pmod {p^2}$ either has no solution or specifies one
residue class of $n$.  Since $H_n\geq x^\delta\gg p^2$, its count
is $O(H_n/p^2)$.  Finally
$\sum_{i,p\in\mathcal P_i}p^{-2}\ll_K\exp(-L^{0.1})$.
The coefficient bounds absorb only a fixed power of $L$, proving
\eqref{det:squareful-error}.

Choose a real smooth dyadic partition with factors $\eta(b_0/Y)$,
where $\eta$ is supported in $[1,4]$ and has bounded derivatives.
There are $O_K(L)$ relevant $Y$, and
\begin{equation}\label{det:Y-range}
 cL^{0.1}\leq\log Y\leq C_KL^{0.2}.
\end{equation}
For a fixed factor of this partition the sum becomes
\[
 \sum_{h\ll X/Y}F(h)q^{\omega(h)}
 \sum_{\substack{mn-1=b_0h\\b_0=\prod_i p_i,\ p_i\in\mathcal P_i}}
       \left(\prod_iV_i^{-1}\right)
       \eta(b_0/Y)\alpha_m\beta_n.
\]
Cauchy's inequality bounds its square by $O(X/Y)$ times the
nonnegative expanded square
\begin{equation}\label{det:cauchy-square}
 \begin{split}
 \mathcal Q_Y={}&\left(\prod_iV_i^{-2}\right)
 \sum_{\substack{a,b,m,n,r,s\\mn-1=ah,\ rs-1=bh\text{ for some }h}}
 \eta(a/Y)\eta(b/Y)
 \alpha_m\overline{\alpha_r}\beta_n\overline{\beta_s}.
 \end{split}
\end{equation}
Both $a$ and $b$ select one prime per group.  In particular, proving
$\mathcal Q_Y\ll XYL^{-D_1}$ for arbitrarily large fixed $D_1$
is sufficient; the required $D_1$ depends on $C,D_*$ but not on $K$.

We may first discard pairs $a,b$ sharing a label.  The elementary
bound used here, uniform for $l\asymp Y$, is
\begin{equation}\label{det:progression-divisors}
 \sum_{h\ll X/Y}\tau(1+lh)^2\ll (X/Y)L^3.
\end{equation}
Indeed, $\tau(n)^2\leq\tau_4(n)$.  In an ordered four-factor
decomposition of an integer at most $O(X)$, deleting a largest
factor leaves a product $d\ll X^{3/4}$.  There are at most
$4\tau_3(d)$ choices for the three retained factors and their
positions.  If $d\mid1+lh$, then $(d,l)=1$, and $h$ lies in one
progression modulo $d$.  Consequently the left side of
\eqref{det:progression-divisors} is at most a constant times
\[
 \sum_{d\ll X^{3/4}}\tau_3(d)\left(\frac{X}{Yd}+1\right)
 \ll \frac XY L^3+X^{3/4}L^2\ll\frac XY L^3.
\]
These divisor estimates also follow from Lemma~\ref{lem:moment-bounds}.
For fixed $a,b$, Cauchy's inequality and
\eqref{det:progression-divisors} bound the absolute $h$-sum of the
two representation counts by $(X/Y)L^{O_C(1)}$.
The number of pairs sharing a prime is
\[
 \ll Y^2\sum_{i,p\in\mathcal P_i}p^{-2}
 \ll_K Y^2\exp(-L^{0.1});
\]
here the number of integers in $[Y,4Y]$ divisible by $p$ is
$O(Y/p)$.  Thus the discarded part is $O_A(XYL^{-A})$.

For the remaining pairs $(a,b)=1$.  Their two equations in
\eqref{det:cauchy-square} are equivalent to
\begin{equation}\label{det:equation}
 t=bmn-ars=b-a.
\end{equation}
In fact, the displayed equality gives $a\mid mn-1$ and
$b\mid rs-1$, and the quotients agree and are positive.

Fix a real compactly supported smooth $\psi$, equal to one on
$[-4,4]$.  For a fixed exponent $A_0>0$, let $\mathfrak M$ be the
union, on $\R/\Z$, of the arcs
\[
 \left|\theta-\frac hk\right|\leq\frac{2L^{A_0}}Y,
 \qquad 1\leq k\leq L^{A_0},\quad (h,k)=1.
\]
They are disjoint for large $x$.  The replacement for
\eqref{det:equation} is
\begin{equation}\label{eq:det-major}
 \mathcal H_{\mathfrak M}(t;a,b)
   =\psi(t/Y)\int_{\mathfrak M}
                  \e\bigl(\theta(t-b+a)\bigr)\,\dd\theta.
\end{equation}
In the ultimate major sum $a,b$ need not be disjoint.  The signed
error in this replacement is controlled by the lift constructed
next.

\subsection{Physical states, good positions, and endpoint vectors}

Use the growing parameters
\begin{equation}\label{det:moment-parameters}
 J=\lfloor L^{0.01}\rfloor,\quad M=J+1,\quad
 R=\lceil L^{0.5}/2\rceil,\quad N=2R.
\end{equation}
An auxiliary list has $J$ primes per group; its product is denoted
by $D$.  Partition all possible $D$ into ordinary dyads
$[d_0,2d_0)$.  There are $O_K(L)$ such dyads and
$\log(2d_0)\leq C_KL^{0.21}$.  Fix one and put
\begin{equation}\label{det:physical-box}
 H=d_0Y,\qquad U=HH_m,\qquad V=H_n,\qquad
 \Omega=[U,16U]\times[V,2V].
\end{equation}

A \emph{physical state} is a primitive integer vector
$P=(P_1,P_2)\in\Omega$ together with an ordered list
$\boldsymbol p_i=(p_{i,1},\ldots,p_{i,M})$ of distinct primes
from $\mathcal P_i$ for each $i$, all dividing $P_1$.
Let $\mathcal H$ be the Hilbert space on these states with measure
\begin{equation}\label{det:state-measure}
 \dd\sigma(P,\boldsymbol p)=\prod_{i=1}^KV_i^{-M}
 \quad\text{times counting measure}.
\end{equation}
Averaging independently over permutations of each list defines an
orthogonal projection $\mathcal S$ on $\mathcal H$.

We now specify the row operation before symmetrization.  Copy slots
$1,\ldots,J$ in every group from source to target; their product
is $D$.  The product of the source's last slots is $b$.  The
target's last slots are new labels with product $a$.  Each new
slot is summed with weight $V_i^{-1}$; target positions are summed
with counting measure.  The labels shared across the edge are
exactly the copied slots.  Thus all unshared labels are distinct
from one another and from every shared label, including across
the two endpoints.  This restriction will be called the
\emph{cross-edge ban}.  Require $D\in[d_0,2d_0)$ and
$t=\det(P,Q)/D\in\Z$, and use the multiplier
\begin{equation}\label{eq:det-kernel}
 \begin{split}
 \mathcal K(P,Q;D,a,b)={}&\frac{d_0}{D}
       \eta(b/Y)\eta(a/Y)\psi(t/Y)\\
 &\times\left\{\1_{t=b-a}
       -\int_{\mathfrak M}\e\bigl(\theta(t-b+a)\bigr)
                    \,\dd\theta\right\}\\
 &\times q^{(\omega(P_1)-KM)/2+(\omega(Q_1)-KM)/2}.
 \end{split}
\end{equation}
Let $\mathcal T$ be this row operation. All sums are finite.
The adjoint reverses the row rule and
conjugates its multiplier.  Indeed, the joint measure of the two
lists in an edge pairing is $\prod_iV_i^{-(M+1)}$ in either
direction: one factor $V_i^{-M}$ comes from
\eqref{det:state-measure}, and one $V_i^{-1}$ from the new slot.

After slot symmetrization, the copied products $D$ range over all choices
obtained by omitting one prime from each full list and lying in the chosen
pad dyad $[d_0,2d_0)$. Each of the $M^K$ omission tuples has weight
$M^{-K}$ under the source symmetrization.

We next restrict the positions at which this row operation is used.
The first restriction will keep the lattice determined by the shared
labels from becoming too thin. The second separates phases associated
with the omission choices; in the repeated-prime argument it will leave
at most one compatible omission tuple.

For a primitive $P$, complete it to an integral matrix of determinant
one, and let $r_P$ be the first coordinate of the second column
divided by $P_1$, considered modulo one.  Changing the completion
adds an integer, so $r_P$ is well defined.  A position with its
full unordered lists is \emph{good} if the following hold.
\begin{enumerate}[label=(\roman*)]
\item\label{det:good-first} For every $D\in[d_0,2d_0)$ obtained by
omitting one label per group, there is no integer
$1\leq l\leq Y^{0.2}$ such that
$\|lDr_P\|_{\R/\Z}\leq Y^{-0.7}$.
\item\label{det:good-second} Let $I$ be any nonempty subset of the
groups and $i_* =\max I$.  Draw one fresh prime in each group
outside $I$, independently according to $\mu_i$, and let $T_f$
be their product.  Except for a set of these draws of probability
at most $\exp(-\tfrac14L^{a_{i_*}})$, the points
$DZ T_f r_P$ corresponding to all distinct integers $DZ$ are
pairwise separated in circle distance by
$100\exp(-L^{a_{i_*}})$.  Here $D$ ranges over the omission
products just described, and $Z$ selects one prime from each
group in $I\setminus\{i_*\}$.
\end{enumerate}
Let $\mathcal G$ be the projection onto good states. The definition
is symmetric in the ordered slots, so $\mathcal G$ commutes with
$\mathcal S$. We use the restricted, symmetrized operator
\[
 \mathcal A=\mathcal G\mathcal S\mathcal T\mathcal S\mathcal G.
\]
The first test is used in the lattice construction leading to
\eqref{mem:lattice-shape}; the second is used to prove the unique
omission choice in the argument following \eqref{mem:anchor-congruence}.

\begin{lemma}\label{det:good-state}
For any fixed lists, the set of $r\in\R/\Z$ failing goodness has
Lebesgue measure at most $\exp(-cL^{0.1})$, for some $c>0$ and
all sufficiently large $x$.
\end{lemma}
\begin{proof}
There are at most $M^K$ omission products.  The union bound for
the first test is at most $2M^KY^{-0.5}$, which has the required
size by \eqref{det:Y-range}.  For the second test fix $I,T_f$.
The number of trial integers $DZ$ is
\[
 M^K\prod_{i\in I\setminus\{i_*\}}|\mathcal P_i|
          =\exp\bigl(o(L^{a_{i_*}})\bigr).
\]
For distinct $DZ,D'Z'$, the integer
$(DZ-D'Z')T_f$ is nonzero.  Multiplication by this integer
preserves uniform measure on the circle.  The measure on which
this pair violates the required separation is therefore at most
$200\exp(-L^{a_{i_*}})$.  Summing over pairs and then averaging
over $T_f$ gives $\exp(-L^{a_{i_*}}+o(L^{a_{i_*}}))$.
Markov's inequality at the threshold
$\exp(-\tfrac14L^{a_{i_*}})$ gives
$\exp(-\tfrac34L^{a_{i_*}}+o(L^{a_{i_*}}))$ for the exceptional
set of $r$.  There are only $2^K-1$ possible $I$.
\end{proof}

We record precisely the vectors that recover the Cauchy square.
At an endpoint require the complete part of $P_1$ supported on
the group primes to be squarefree, with exactly $M$ primes from
each group.  Denote that part by $G(P_1)$.  Put
\begin{equation}\label{det:endpoints}
 \overline{f(P)}=\alpha_{P_1/G(P_1)}\overline{\beta_{P_2}},
 \qquad
 g(P)=\overline{\alpha_{P_1/G(P_1)}}\beta_{P_2},
\end{equation}
and set these vectors to zero when the stated conditions or the
coefficient supports fail.  The vectors are constant on lists.
In an edge contributing to their pairing, the coordinates are
\begin{equation}\label{det:lift}
 P=(Dbm,s),\qquad Q=(Dar,n),\qquad
 \det(P,Q)/D=bmn-ars.
\end{equation}
No new restriction is imposed on the original coefficients:
every $m,r,n,s$ with nonzero coefficient is rough, and every
group prime is below $W$.  The damping in
\eqref{eq:det-kernel} equals one at these endpoints.

Their norms have the uniform bounds
\begin{equation}\label{det:endpoint-norms}
 \|f\|,\|g\|\ll(UV)^{1/2}L^{O_C(1)},\qquad
 \sup|f|+\sup|g|\ll L^{O_C(1)}.
\end{equation}
For example, at a fixed ordered list with product $G$, the number
of possible positions is $O(UV/G)$, since $G=x^{o(1)}\ll U$.
Summing its reciprocal product with the state normalization gives
\[
 \prod_i V_i^{-M}
 \sum_{\substack{p_{i,1},\ldots,p_{i,M}\in\mathcal P_i\\
                  \text{distinct in each group}}}
             \prod_{i,j}\frac1{p_{i,j}}\leq1.
\]
This proves \eqref{det:endpoint-norms} without a logarithmic
exponent depending on $K$ or $J$.

Conversely, positions in \eqref{det:lift} on the support of
\eqref{eq:det-kernel} are automatically primitive.  A prime
dividing both $Dbm$ and $s$ cannot divide $Db$, by roughness.
It would therefore divide $m,s$ and be larger than $W$.  It then
divides $t=bmn-ars$.  The cutoff gives $|t|\ll Y<W$, so $t=0$.
But $bmn=ars$ is impossible: any prime in $b$ divides none of
$a,r,s$.  The same proof applies to $Q$.  The ranges in
\eqref{det:lift} lie in \eqref{det:physical-box} because
$D\in[d_0,2d_0)$ and $a,b\in[Y,4Y]$.

\subsection{From the moment to endpoint pairings}

For a primitive $P\in\Omega$, let $\mathbf u_P$ be one on every
list at $P$ and zero at other positions.  These vectors are not
normalized.  The estimate proved in
Proposition~\ref{prop:minor-moment} is, for any prescribed $E_0>0$,
\begin{equation}\label{eq:det-moment}
 \frac1{UV}\sum_{P\in\Omega\text{ primitive}}
       \langle\mathbf u_P,(\mathcal A\mathcal A^*)^R
                     \mathbf u_P\rangle
             \leq L^{-E_0N}.
\end{equation}
First $A_0$ and then $K$ are chosen sufficiently large; thresholds
in absolute errors may subsequently depend on the fixed group
exponents.  We prove here the consequence
\begin{equation}\label{eq:det-pairing}
 |\langle f,\mathcal A g\rangle|
             \ll UVL^{-E_0+O_C(1)}.
\end{equation}

Partition positions into half-open intervals of the slope $P_2/P_1$
of length $H/U^2$.  A nonzero edge has
$|\det(P,Q)|\ll H$, so only a bounded number of neighboring
intervals can interact.  Each interval contains $O(H)$ primitive
positions.  To see this, choose a primitive pivot in it.  Its
determinant with any other position in the interval is an integer
of size $O(H)$.  With that determinant fixed, all integer
solutions differ by integral multiples of the pivot; the first
coordinate ranges over $[U,16U]$, permitting $O(1)$ multiples.

Let $f_j$ be $f$ restricted to one interval, and $g_j$ the
restriction of $g$ to its interacting neighbors.  Set
$B_0=\mathcal A\mathcal A^*$ and
\[
 d_j=\sum_{P\text{ in interval }j}
              \langle\mathbf u_P,B_0^R\mathbf u_P\rangle.
\]
The matrix $C^{(j)}_{PQ}=\langle\mathbf u_P,B_0^R\mathbf u_Q\rangle$
is positive semidefinite on the ordinary coefficient space.
Its largest eigenvalue is at most its trace $d_j$, irrespective
of the norms or orthogonality of the testing vectors.  Since
$f_j=\sum_P f(P)\mathbf u_P$, this gives
\begin{equation}\label{det:compression}
 \langle f_j,B_0^Rf_j\rangle
      \leq d_j\sum_P|f(P)|^2
      \ll H\sup|f|^2d_j.
\end{equation}
Spectral H\"older for the positive operator $B_0$, followed by
Cauchy's inequality, now gives
\[
 |\langle f_j,\mathcal A g_j\rangle|
 \leq \|g_j\|\,\|f_j\|^{1-1/R}
             \bigl(CH\sup|f|^2d_j\bigr)^{1/(2R)}.
\]
H\"older's inequality in $j$, with exponents
$2,2R/(R-1),2R$, and bounded overlap of the $g_j$ imply
\[
 |\langle f,\mathcal A g\rangle|
 \ll\|g\|\,\|f\|^{1-1/R}
       (CH\sup|f|^2)^{1/(2R)}
       \left(\sum_jd_j\right)^{1/(2R)}.
\]
By \eqref{det:Y-range} and \eqref{det:physical-box},
$\log H=O_K(L^{0.21})$, so $H^{1/N}=O(1)$.
Equations~\eqref{eq:det-moment} and
\eqref{det:endpoint-norms} prove \eqref{eq:det-pairing}.
This argument uses the unnormalized trace in
\eqref{eq:det-moment}; no normalization by the number of lists
has been introduced.

\subsection{Uniform residues at a primitive root}

The root average in \eqref{eq:det-moment} will be replaced by
independent projective residue lines.  We give an elementary
count that works even when the two coordinate scales are
comparable.  For a primitive $(u,v)$ in $\Omega$ there is a
unique completion
\begin{equation}\label{det:root-matrix}
 g=\begin{pmatrix}u&c\\v&d\end{pmatrix},\qquad
 ud-vc=1,\qquad 0\leq c<u.
 \quad\text{Put }r=c/u.
\end{equation}

\begin{lemma}[Root residues]\label{lem:root-residue}
There is $c_\delta>0$ such that the following holds.  Let $S$ be
squarefree with $S\leq\exp(C_KL^{0.98})$, and prescribe
$g_0\in\mathrm{SL}_2(\Z/S\Z)$.  If $I_1,I_2,I_3$ are intervals
in $[1,16],[1,2],[0,1]$, respectively, the number of primitive
roots satisfying
\[
 u/U\in I_1,\qquad v/V\in I_2,\qquad r\in I_3,
 \qquad g\equiv g_0\pmod S
\]
is
\begin{equation}\label{det:root-residue-count}
 \frac{UV\,|I_1||I_2||I_3|}
      {\zeta(2)|\mathrm{SL}_2(\Z/S\Z)|}
             +O(UVx^{-c_\delta}).
\end{equation}
The error is uniform in the intervals and in $g_0$.
\end{lemma}

\begin{proof}
We first prove the elementary exponential-sum estimate needed for
the count.  For $m\geq1$ put
\[
 K_m(h,k)=\sum_{y\in(\Z/m\Z)^\times}
                   \e\left(\frac{hy+ky^{-1}}m\right).
\]
Orthogonality shows that
$\sum_{h,k\bmod m}|K_m(h,k)|^4=m^2T_m$, where $T_m$ counts
unit quadruples $(y_1,y_2,z_1,z_2)$ with equal sums and equal
inverse sums.  The first relation determines
$z_2=y_1+y_2-z_1$.  After multiplying the inverse-sum relation
by the product of the four units, one obtains
\[
 m\mid(y_1+y_2)(z_1-y_1)(z_1-y_2).
\]
The map from the three free residues to these three linear forms
has determinant of absolute value two, and its kernel modulo
$m$ has size at most two.  Distribute, prime by prime, the
valuation of $m$ among the three factors.  There are
$\tau_3(m)$ distributions $d_1d_2d_3=m$; for each, the number
of triples of forms with the prescribed divisibilities is
$m^3/(d_1d_2d_3)=m^2$.  Consequently
$T_m\leq2m^2\tau_3(m)$.

The value of $K_m(h,k)$ is constant on the unit-scaling orbit
$(h,k)\mapsto(hs,ks^{-1})$.  If $g=(h,k,m)$, its stabilizer
consists of the units $s\equiv1\pmod{m/g}$, so the orbit has
$\varphi(m/g)\geq(m/g)m^{-o(1)}$ elements.  Dividing the fourth
moment bound by this orbit size gives
\begin{equation}\label{det:weak-kloosterman}
 |K_m(h,k)|\ll m^{3/4+o(1)}(h,k,m)^{1/4}.
\end{equation}
Only the elementary bounds $\tau_3(m)=m^{o(1)}$ and
$\varphi(m)\geq m^{1-o(1)}$ enter here; both follow by separating
the finitely many small primes from the remaining prime factors.

Suppose first $U\leq V$.  Write the entries of the prescribed
matrix as $u_0,c_0,v_0,d_0'$; the prime on $d_0'$ distinguishes
it from the pad scale.  Fix $u\equiv u_0\pmod S$.  Modulo
$m_0=uS$, the required pairs $(c,v)$ satisfy
\begin{equation}\label{det:fixed-u-congruences}
 c\equiv c_0\pmod S,\quad v\equiv v_0\pmod S,
 \quad cv\equiv ud_0'-1\pmod{uS}.
\end{equation}
There are exactly
\begin{equation}\label{det:fixed-u-mass}
 M(u)=u\prod_{\substack{p\mid u\\p\nmid S}}(1-p^{-1})
\end{equation}
such pairs.  Indeed, the admissible $c$ must be coprime to $u$.
At primes common to $u,S$ this is forced by
$u_0d_0'-c_0v_0=1$ modulo $S$; at other primes the indicated
Euler factors count the admissible $c$ in its progression.
For any such $c$, writing $v=v_0+Sj$ reduces the last congruence
to a linear congruence modulo $u$ with invertible coefficient
$c$, and hence determines exactly one $j\pmod u$.

These $M(u)$ pairs have discrepancy
\begin{equation}\label{det:modular-discrepancy}
 O(u^{0.99}S^2)
\end{equation}
for rectangles in the torus $(c/(uS),v/(uS))$.  Here are details
of the uniformity in $S$.  Put
$S_1=\prod_{p\mid(S,u)}p$ and $S_2=S/S_1$.  The Chinese
remainder theorem fixes $(c,v)$ modulo $S_2$ and leaves an
inverse graph modulo $m_1=uS_1$.  Write $a'=ud_0'-1$ and
$\gamma=S_2^{-1}\pmod {m_1}$; both are units modulo $m_1$.
On the inverse graph $cv=a'$ the condition $c\equiv c_0\pmod
{S_1}$ already forces $v\equiv v_0\pmod {S_1}$, since $c_0$
is a unit there.  Additive orthogonality thus gives the Fourier
sum, up to a constant of absolute value one, as
\begin{equation}\label{det:CRT-Fourier}
 \frac1{S_1}\sum_{j\bmod S_1}\e(-jc_0/S_1)
          K_{m_1}(\gamma h+ju,\gamma k a').
\end{equation}
For $0<\max(|h|,|k|)\leq u^{0.02}$,
\[
 (\gamma h+ju,\gamma ka',m_1)
      \leq S_1(h,k,u)\leq S_1u^{0.02}.
\]
Using the exponent $3/4+1/200$ in
\eqref{det:weak-kloosterman}, each nonzero sum
\eqref{det:CRT-Fourier} is therefore $O(u^{0.76}S^2)$.

For completeness, sandwich each interval indicator between
smooth upper and lower functions obtained by enlarging or
contracting the interval by $2u^{-0.01}$ and convolving with
a probability bump of width $u^{-0.01}$.  Their zeroth
coefficients differ from the interval length by $O(u^{-0.01})$.
Their Fourier coefficients are bounded.  There are $O(u^{0.04})$
retained frequency pairs, whose total cost is $O(u^{0.80}S^2)$.
For explicit control of the tail, if $\Delta=u^{-0.01}$ and
$T_1=u^{0.02}$, ten integrations by parts give a double
Fourier tail outside $\max(|h|,|k|)\leq T_1$ bounded by
$O(\Delta^{-11}T_1^{-9})=O(u^{-0.07})$.  The trivial bound
$M(u)\leq u$ makes its contribution $O(u^{0.93})$.
The error from the zeroth coefficient is $O(u^{0.99})$.
This proves \eqref{det:modular-discrepancy}.

The $c$ interval has length $u|I_3|$, and the $v$ interval has
length $V|I_2|$.  If the latter passes through several periods
of length $uS$, apply the rectangle count to each; the number
of pieces is $O(V/u+1)$.  Equations
\eqref{det:fixed-u-mass}--\eqref{det:modular-discrepancy} give
the fixed-$u$ main term
\[
 \frac{V|I_2||I_3|}{S^2}
           \prod_{\substack{p\mid u\\p\nmid S}}(1-p^{-1})
\]
and error $O((V/u+1)u^{0.99}S^2)$.
For $u\asymp U$ the sum of these errors is
$O(UV U^{-0.01}S^2)$.

It remains to average the Euler factor in the progression of
$u$.  Expanding it as
$\sum_{l\mid u,(l,S)=1}\mu(l)/l$ gives
\[
 \sum_{\substack{u/U\in I_1\\u\equiv u_0\ (S)}}
 \prod_{\substack{p\mid u\\p\nmid S}}(1-p^{-1})
 =\frac{U|I_1|}{S}
       \sum_{(l,S)=1}\frac{\mu(l)}{l^2}+O(\log(2U)).
\]
The tail beyond $16U$ contributes $O(1/S)$ and is covered by
the error.  The series equals
$\zeta(2)^{-1}\prod_{p\mid S}(1-p^{-2})^{-1}$.
Counting first columns and their completions over each field
shows
\begin{equation}\label{det:SL-order}
 |\mathrm{SL}_2(\Z/S\Z)|
       =S^3\prod_{p\mid S}(1-p^{-2}).
\end{equation}
This proves the stated main term.  Since $U\geq x^\delta$ and
$S^2=x^{o(1)}$, all the errors admit a saving $x^{-c_\delta}$;
for example, $c_\delta=\delta/1000$ is sufficient.

If $V<U$, fix $v$ instead and apply the same argument to the
congruence $ud\equiv1+vc_0\pmod{vS}$, with the roles of the
coordinates reversed and determinant sign reversed.  Use
$d/v$ as the third coordinate.  For $u,v>1$ both $c/u$ and
$d/v$ belong to $[0,1)$, and
\[
 \frac dv-\frac cu=\frac1{uv}.
\]
Enlarging and contracting the prescribed interval by
$O((UV)^{-1})$ changes its main mass by $O(1)$ and its count
by the already available discrepancy.  The same formula follows.
\end{proof}

\subsection{Independent residue lines and the archimedean error budget}

Expand \eqref{eq:det-moment} as a path with $N$ edges, returning
to its initial position but with no return condition on the
lists.  In the basis \eqref{det:root-matrix}, each position is
$gz$ with primitive $z\in\Z^2$, and the initial and final
vectors are $e_1=(1,0)$.  Put $\tau(z)=z_1+rz_2$.  The cutoff
and the physical box imply
\begin{equation}\label{det:tube}
 |z_2|\ll NH,\qquad |z_1|\ll NH,\qquad \tau(z)\asymp1.
\end{equation}
Indeed, one edge changes slope by $O(H/U^2)$; summing at most
$N$ changes gives $z_2=\det((u,v),gz)=O(NH)$.
Also $(gz)_1=u\tau(z)\asymp U$, which proves the remaining
claims.  For a group prime $p$, the condition $p\mid(gz)_1$
means that the projective line $[z]_p$ equals the kernel line
of the first row of $g$ modulo $p$.

\begin{lemma}[Replacement of the root average]\label{det:root-replacement}
In the normalized path expansion of \eqref{eq:det-moment}, one
may replace the root average by
\[
 \frac1{\zeta(2)}\int_{[1,16]\times[1,2]\times[0,1]}
                  \dd(u/U)\,\dd(v/V)\,\dd r
\]
and by independent uniformly distributed lines in
$\mathbb P^1(\mathbb F_p)$ for all the group primes.  The real
matrix at an archimedean root is
\[
 g(u,v,r)=\begin{pmatrix}u&ur\\v&vr+1/u\end{pmatrix}.
\]
All box and good-state conditions are retained.  The total error
is $O_A(L^{-AN})$ for every fixed $A>0$.  The same assertion,
with polynomial logarithmic coefficient bounds, holds for one
edge and for failure of either endpoint's goodness.
\end{lemma}

\begin{proof}
We specify both the combinatorial and the archimedean budgets.
There are $\exp(O_K(L^{0.73}))$ choices for paths $z$ satisfying
the size bounds in \eqref{det:tube}, labels of all visits, and
slot permutations.  For example, the logarithm of the number
of position paths is $O(N\log(NH))=O_K(L^{0.71})$;
the label choices have logarithm at most
$O_K(NML^{0.2})=O_K(L^{0.71})$; and the permutations cost
$O_K(NM\log M)$.  These also bound the total absolute
coefficients after dropping congruences and bounded damping.
The product of the primes active anywhere on a path has
logarithm $O_K(L^{0.72})$.

Keep the exact factors of the primes active at some visit.
At the other primes the joint damping can be written $1-X_p$,
where $0\leq X_p\leq1$.  Here the damping exponent at an
endpoint is $q_j=\sqrt q$ and at an interior visit is $q_j=q$.
In the independent-line model,
\[
 \sum_p\E X_p
 \leq\sum_p\sum_{j=0}^N\frac{1-q_j}{p+1}
 \ll_K N+1.
\]
Upper and lower Bonferroni polynomials of consecutive degrees
near $T=\lceil L^{0.76}\rceil$ sandwich
$\prod_p(1-X_p)$ pointwise.  Their expected difference is at
most
\[
 \frac{(C_K(N+1))^T}{T!}
       =\exp\bigl(-\Omega(L^{0.76}\log L)\bigr).
\]
The product inequality is ordinary inclusion--exclusion, applied
to numbers in $[0,1]$; it requires no independence for the
pointwise sandwich.  Independence between different residue
lines is used only to estimate its gap.  This gap absorbs the
$\exp(O_K(L^{0.73}))$ path and coefficient budget.

Each monomial in the truncated expressions requires residues at
the active primes and at most $T$ further primes.  Its squarefree
modulus satisfies
$\log S=O_K(L^{0.96})$, and hence lies in the range of
Lemma~\ref{lem:root-residue}.  The total number of monomials,
prime choices, and residue matrices is $\exp(o(L))$.
The uniform distribution on $\mathrm{SL}_2(\Z/S\Z)$ factors
over the primes by the Chinese remainder theorem.  At each
prime its first row is uniform among nonzero rows, and its
kernel is uniform among the $p+1$ projective lines.  This is
exactly the asserted independent model, with primitive-root
density $1/\zeta(2)$.

We next make the archimedean conditions compatible with interval
counts.  Choose an integral complement $w$ to each primitive
$z$, with $\det(z,w)=1$ and $|w|\ll1+|z|$.  The ratio for the
position $gz$ is
\begin{equation}\label{det:transported-ratio}
 r_{gz}=\frac{\tau(w)}{\tau(z)}\pmod1.
\end{equation}
On $\tau(z)\asymp1$ its derivative in $r$ is
$1/\tau(z)^2$.  The good tests are therefore constant except
at $\exp(o(L))$ values of $r$: enumerate every product and
fresh draw in the tests, every pair of distinct integers $DZ$,
and the integer translates in each circle inequality.
All their numbers and sizes are $\exp(o(L))$.
The exceptional-probability condition in the second good test
is also constant between these endpoints, because it is a
finite weighted sum of such indicators.

The physical coordinates are
\begin{equation}\label{det:transported-position}
 (gz)_1=u\tau(z),\qquad
 (gz)_2=v\tau(z)+z_2/u.
\end{equation}
The normalized box faces consequently have derivatives bounded
by $\exp(o(L))$.  Near a first-coordinate face the derivative
in $u/U$ is bounded below, and near a second-coordinate face
the derivative in $v/V$ is bounded below, because
$\tau(z)\asymp1$.  One may first exclude $\tau(z)$ near zero,
which is incompatible with the first-coordinate box condition.
Thus the union of cubes meeting any face or good-test boundary
has volume at most $x^{-\epsilon_\delta+o(1)}$ on a grid of
side $x^{-\epsilon_\delta}$ in $(u/U,v/V,r)$.

Choose the grid exponent only after the saving in
Lemma~\ref{lem:root-residue}; concretely take
$\epsilon_\delta=c_\delta/8$.  The grid has
$O(x^{3\epsilon_\delta})$ cubes, not a subpower number.
Summing all root-count errors, including the
$\exp(o(L))$ discrete choices, costs at most
\begin{equation}\label{det:grid-discrepancy}
 UVx^{-c_\delta+3\epsilon_\delta+o(1)}
       =UVx^{-5c_\delta/8+o(1)}.
\end{equation}
The main mass of boundary cubes is at most
\begin{equation}\label{det:grid-boundary}
 UVx^{-\epsilon_\delta+o(1)}
       =UVx^{-c_\delta/8+o(1)}.
\end{equation}
Their discrepancy is already included in
\eqref{det:grid-discrepancy}.  On each remaining cube all
archimedean indicator conditions are fixed.  The edge factors
in \eqref{eq:det-kernel} depend, after fixing labels, on
$\det(z,z')/D$ and hence do not vary with the root; the
finite residue factors have already been handled.
This proves the replacement with a fixed-power error after
normalization.  Such an error is $O_A(L^{-AN})$, since
$N\asymp L^{0.5}$.  The one-edge versions have smaller
enumeration budgets and use the same proof.
\end{proof}

It remains to bound the signed independent-line path expansion.
The next section keeps the precise correlations caused by a
prime that leaves an active list and later returns, proves
\eqref{eq:det-moment}, and then removes the good-state
projections and the auxiliary pads.  This yields the
replacement of \eqref{det:cauchy-square} by
\eqref{eq:det-major} with error $O_A(XYL^{-A})$.
